\documentclass[11pt]{article}
\usepackage[margin=1in]{geometry}
\usepackage[T1]{fontenc}
\usepackage[utf8]{inputenc}
\usepackage{lmodern,microtype,amsmath,amssymb,amsthm,mathtools}
\usepackage[colorlinks=true,linkcolor=blue,citecolor=blue,urlcolor=blue]{hyperref}
\hypersetup{pdftitle={Small circuits can have exponentially large prior complexity below the mean-field scaling},pdfauthor={Dmitry Vaintrob; Kaarel Hanni; OpenAI Codex}}
\setlength{\emergencystretch}{2em}
\newtheorem{theorem}{Theorem}[section]
\newtheorem{lemma}[theorem]{Lemma}
\newtheorem{proposition}[theorem]{Proposition}
\newtheorem{corollary}[theorem]{Corollary}
\theoremstyle{remark}
\newtheorem{remark}[theorem]{Remark}
\DeclareMathOperator{\sign}{sign}
\DeclareMathOperator{\Var}{Var}
\newcommand{\R}{\mathbb R}
\newcommand{\E}{\mathbb E}
\newcommand{\Prb}{\mathbb P}
\title{Small circuits can have exponentially large prior complexity\\below the mean-field scaling}
\author{Dmitry Vaintrob \and Kaarel H\"anni \and OpenAI Codex}
\date{October 15, 2026}
\begin{document}
\maketitle
\begin{abstract}
We study a fully connected Bayesian $\tanh$ network on the Boolean
cube, with Gaussian layer variances on the scaling line indexed by
$\gamma$, and a prior conditioned to have a fixed output margin at
every input. For every fixed $\gamma<1$, we show that a small circuit
need not have polynomial prior complexity. Full parity on $d$ bits
has a depth-$L$ implementation with width $d+1$ and weights
$O(d\log d)$, for every $L\ge2$. Nevertheless, for $\gamma\le0$ its
limiting normalized prior probability is at most $2^{-2^d}$. For
$0<\gamma<1$, its negative log prior probability, divided by
$N^\gamma$, is at least $a^2(2^d-O_L(d))/2$ in the width limit,
where $a$ is the margin. The proof uses the exact spectrum of the
saturated neural-network kernel and a general Gaussian inequality:
a Boolean minimum-eigenvalue eigenvector has probability at most
$2^{-M}$ among sign patterns, even after conditioning on any common
global margin. These are prior-mass separations. They do not by
themselves establish exponential sample complexity for the resulting
Bayesian classifier.
\end{abstract}

\section{Model and the prior-complexity question}

Let $\mathcal X=\{-1,1\}^d$, $M=2^d$, and $L\ge2$. The input is
augmented by a constant coordinate, $h^0(x)=(1,x)$. A network with
$L-1$ hidden layers of common width $N$ has
\[
 h^s(x)=\tanh(W^s h^{s-1}(x)),\qquad
 z_N(x)=W^L h^{L-1}(x).
\]
The readout is linear, and there are no further bias coordinates.
All weights are independent centered Gaussians with variances
\begin{equation}\label{eq:scaling}
 \Var(W^1_{ij})=N^{1-\gamma}/d,\qquad
 \Var(W^s_{ij})=N^{-\gamma}\quad(1<s<L),\qquad
 \Var(W^L_i)=N^{-\gamma-1}.
\end{equation}
Fix $a>0$. Write $Z_N=(z_N(x))_{x\in\mathcal X}$, and define
\[
 A_a=\{|Z_N(x)|>a\text{ for every }x\},\qquad
 A_f=\{f(x)Z_N(x)>a\text{ for every }x\}.
\]
The prior on Boolean functions is
\begin{equation}\label{eq:prior}
 \pi_{N,a}(f)=\Prb(A_f)/\Prb(A_a).
\end{equation}
Thus the all-input margin restriction is applied before observing
any training data. It is not a margin imposed only on the sample.

Let $C_L(f)$ be the smallest integer $C$ such that a depth-$L$
network of width at most $C$, with all weights of magnitude at most
$C$, satisfies $f(x)z(x)>1/2$ everywhere. This combines width and
weight size in one circuit-complexity parameter.

The natural prior complexity changes with the overall output
scale. For the present paper define
\begin{equation}\label{eq:complexity}
 \mathcal C_{\gamma,a}(f)=
 \begin{cases}
 \displaystyle\lim_{N\to\infty}-\log\pi_{N,a}(f),&\gamma\le0,\\[4pt]
 \displaystyle\lim_{N\to\infty}-N^{-\gamma}\log\pi_{N,a}(f),
          &0<\gamma<1.
 \end{cases}
\end{equation}
We prove below that these limits exist. The normalization for
positive $\gamma$ removes the Gaussian small-noise speed. Without
it, negative log probabilities would diverge even for many simple
functions.

All width limits take $N\to\infty$ at fixed $d,L,\gamma,a$.
Statements about growth in $d$ concern the resulting limiting
quantities. No uniform exchange of the width and dimension limits
is made.

\begin{theorem}[A circuit-prior separation throughout $\gamma<1$]
\label{thm:main}
Let $f_d(x)=\prod_{j=1}^d x_j$ be full parity, and put
$c_L=(2/\pi)^{L-1}$. For every $L\ge2$,
\begin{equation}\label{eq:circuit-bound}
 C_L(f_d)\le \left\lceil(d+1)\log(16(d+1))\right\rceil.
\end{equation}
For every $\gamma\le0$,
\begin{equation}\label{eq:nonpositive-bound}
 \mathcal C_{\gamma,a}(f_d)\ge (\log2)2^d.
\end{equation}
For every $0<\gamma<1$,
\begin{equation}\label{eq:positive-bound}
 \mathcal C_{\gamma,a}(f_d)
 \ge \frac{a^2}{2}\left(2^d-\frac{d+1}{c_L}\right).
\end{equation}
Consequently, at each fixed depth and fixed subcritical exponent,
there is no polynomial in $d$ and $C_L(f)$ that upper-bounds
$\mathcal C_{\gamma,a}(f)$ for every Boolean target $f$.
\end{theorem}

This theorem answers a prior-mass question: throughout this scaling
regime, polynomial-size circuits need not receive polynomial
description-length prior mass, on the natural scale of the prior.
The distinction between this statement and a posterior learning
lower bound is discussed in Section~\ref{sec:learning}.

\section{The saturated kernel and its spectrum}

Put
\begin{equation}\label{eq:kernel}
 H(t)=\frac2\pi\arcsin t,\qquad
 F=H^{\circ(L-1)},\qquad
 K(x,y)=F\!\left(\frac{1+x\cdot y}{d+1}\right).
\end{equation}
For $U\subseteq\{1,\ldots,d\}$, write
$\chi_U(x)=\prod_{j\in U}x_j$. Translation symmetry on the Boolean
cube makes these parity functions eigenvectors of $K$. Let
$\lambda_r$ denote the eigenvalue of the normalized operator
$M^{-1}K$ on the degree-$r$ parities.

\begin{lemma}[Paired ordering of the eigenvalues]\label{lem:spectrum}
The matrix $K$ is positive definite, and
\begin{equation}\label{eq:paired}
 \lambda_0=\lambda_1>\lambda_2=\lambda_3>
 \lambda_4=\lambda_5>\cdots>0,
\end{equation}
with a final unpaired $\lambda_d$ if $d$ is even. Moreover,
\begin{equation}\label{eq:spectral-bounds}
 \lambda_d\le2^{-d},\qquad
 \lambda_0\ge\frac{c_L}{d+1}.
\end{equation}
\end{lemma}

\begin{proof}
Write $D=d+1$. The arcsine power series has positive coefficients
at every odd degree, zero even coefficients, and sum one at $t=1$.
These properties persist under composition. Thus
$F(t)=\sum_{m\text{ odd}}c_mt^m$, where every $c_m>0$,
$\sum_m c_m=1$, and $c_1=c_L$.

First use the full $D$-dimensional cube and the kernel
$F(u\cdot v/D)$. The degree-$r$ Fourier coefficient of
$(u\cdot v/D)^m$ equals $p_m(r)$, the probability that the simple
coordinate-flip walk reaches a specified vertex at distance $r$
after $m$ steps. Indeed, a term in the power expansion survives
Fourier averaging exactly when the coordinates chosen an odd
number of times form that specified support.

For $0\le r\le D-2$, the walk recursion gives
\begin{align*}
 D[p_{m+1}(r)-p_{m+1}(r+2)]
 ={}&r[p_m(r-1)-p_m(r+1)]\\
 &+(D-r-2)[p_m(r+1)-p_m(r+3)].
\end{align*}
Terms with zero coefficients at the boundary are omitted.
Starting from a point mass at time zero, induction proves
$p_m(r)\ge p_m(r+2)$. The full-cube eigenvalues
$\mu_r=\sum_{m\text{ odd}}c_mp_m(r)$ vanish for even $r$ and
strictly decrease through the odd values of $r$: for odd $r$,
the term $m=r$ is positive at distance $r$ and zero at distance
$r+2$. Each odd $\mu_r$ is positive.

Restricting the distinguished coordinate to $+1$ gives
$\lambda_r=\mu_r+\mu_{r+1}$. This proves
\eqref{eq:paired} and positive definiteness. Since $K(x,x)=1$,
\[
 \sum_{r=0}^d\binom dr\lambda_r=1.
\]
Its smallest eigenvalue therefore satisfies
$\lambda_d\le 1/M$. Finally, the linear term of $F$ contributes
$c_L/(d+1)$ to $\lambda_0$, and all other contributions are
nonnegative. This proves the remaining bound.
\end{proof}

\section{A Gaussian sign inequality that survives global conditioning}

\begin{lemma}[Minimum-eigenvalue signs have at most uniform mass]
\label{lem:gaussian-sign}
Let $\Sigma$ be a positive-definite $M\times M$ covariance matrix,
and let $f\in\{-1,1\}^M$ be an eigenvector for its smallest matrix
eigenvalue $v>0$. If $G\sim\mathcal N(0,\Sigma)$, then for every
$a\ge0$,
\begin{equation}\label{eq:raw-gaussian-bound}
 \Prb\{f_iG_i>a\text{ for every }i\}
 \le \Phi(-a/\sqrt v)^M,
\end{equation}
and
\begin{equation}\label{eq:conditioned-gaussian-bound}
 \Prb\{\sign G=f\mid |G_i|>a\text{ for every }i\}
 \le 2^{-M}.
\end{equation}
Here $\Phi$ is the standard normal distribution function.
\end{lemma}

\begin{proof}
Conjugate the covariance by $\operatorname{diag}(f)$, so the target
sign is the all-ones vector $\mathbf1$ and this vector has
eigenvalue $v$. Write the conjugated Gaussian as
\[
 G'=\sqrt v\,\xi+Y,
\]
where $\xi$ has independent standard normal coordinates and $Y$
is independent with covariance $\Sigma'-vI$. This covariance is
positive semidefinite and kills $\mathbf1$. Consequently
$\sum_iY_i=0$ almost surely. Set $\alpha=\Phi(-a/\sqrt v)$.

The function $\log\Phi$ is concave. Conditional on $Y$, Jensen's
inequality therefore gives
\[
 \Prb\{G'_i>a\ \forall i\mid Y\}
 =\prod_{i=1}^M\Phi\!\left(\frac{Y_i-a}{\sqrt v}\right)
 \le \alpha^M.
\]
Integrating proves~\eqref{eq:raw-gaussian-bound}. To verify the
concavity directly, put $r(t)=\phi(t)/\Phi(t)$. Then
$(\log\Phi)''(t)=-r(t)(t+r(t))<0$. For $t<0$ the last factor is
positive because
$\phi(t)=\int_{-\infty}^t(-s)\phi(s)\,ds>-t\Phi(t)$;
for $t\ge0$ it is immediate.

A one-dimensional Gaussian with variance $v$ is least likely to
lie outside $[-a,a]$ when its mean is zero. Indeed the outside
probability is even in the mean, and its derivative is nonnegative
for nonnegative means. Thus, conditional on $Y$,
\begin{align*}
 \Prb\{|G'_i|>a\ \forall i\mid Y\}
 &=\prod_{i=1}^M\left[
   \Phi\!\left(\frac{Y_i-a}{\sqrt v}\right)
  +\Phi\!\left(\frac{-Y_i-a}{\sqrt v}\right)\right]\\
 &\ge (2\alpha)^M.
\end{align*}
Integrating and dividing the numerator and denominator bounds
gives~\eqref{eq:conditioned-gaussian-bound}. Conjugation preserves
the global absolute-margin event, completing the proof.
\end{proof}

The lemma does not require a quantitative lower bound on $v$ or
an estimate of the global margin probability. In particular it
does not lose a large factor when the conditioning event is rare.

\section{Width limits and the prior-rate calculation}

We include the limiting argument to specify exactly how the
Gaussian calculation applies to the finite networks.

\begin{proposition}[Kernel freezing and output limits]
\label{prop:kernel-limit}
Fix $\gamma<1$. Let $Q_N$ be the empirical covariance of the final
hidden features on the full cube. For every $\epsilon>0$, there
are $c,C>0$ and $N_0$ such that
\begin{equation}\label{eq:freezing}
 \Prb\{\|Q_N-K\|_{\mathrm{op}}>\epsilon\}\le Ce^{-cN}
 \qquad(N\ge N_0).
\end{equation}
Conditional on the hidden layers,
\begin{equation}\label{eq:output-covariance}
 Z_N\sim\mathcal N(0,N^{-\gamma}Q_N).
\end{equation}
For $\gamma=0$, the globally margin-conditioned sign prior
converges to that of $G\sim\mathcal N(0,K)$ with margin $a$.
For $\gamma<0$, it converges to the unconditioned sign law of $G$.
For $0<\gamma<1$, $Z_N$ has a large-deviation principle at speed
$N^\gamma$ with rate $I(z)=z^TK^{-1}z/2$.
\end{proposition}

\begin{proof}
Put $b_N=N^{(1-\gamma)/2}\to\infty$. Conditional on a hidden
layer, the next preactivation rows are independent Gaussian
vectors with covariance $b_N^2$ times that layer's empirical
covariance. The first layer has the deterministic covariance
$b_N^2(1+x\cdot y)/d$, whose correlation matrix is
$(1+x\cdot y)/(d+1)$.

If $U\sim\mathcal N(0,q)$ with $q$ bounded below, direct integration
gives
\[
 \E|\tanh(bU)-\sign U|
 \le \frac{2\log2}{b\sqrt{2\pi q}}.
\]
The correlation of the signs of a centered Gaussian pair with
correlation $r$ is $H(r)$. Thus conditional expected products of
hidden activations approach the normalized arcsine transform,
uniformly in a sufficiently small neighborhood of each limiting
covariance. All distinct-input correlations stay strictly between
$-1$ and $1$, and diagonal variances stay bounded below. The
off-diagonal normalized arcsine map is consequently locally
Lipschitz there; the limiting diagonal is identically one.

For each pair of inputs the $N$ conditional row products lie in
$[-1,1]$ and are independent. Hoeffding's inequality and a union
bound over the fixed $M^2$ pairs give exponential concentration
at speed $N$. Induction over the fixed number of layers, absorbing
the deterministic $O(b_N^{-1})$ error, proves~\eqref{eq:freezing}.
The independent Gaussian readout gives~\eqref{eq:output-covariance}.

It follows that $N^{\gamma/2}Z_N$ converges weakly to the
nondegenerate Gaussian $G$. At $\gamma=0$, all margin orthants and
their unions have boundary of Gaussian measure zero and positive
probability. Their conditional probabilities therefore converge.
At $\gamma<0$, the normalized margin is
$aN^{\gamma/2}\to0$. The probability that any coordinate lies in
this vanishing interval tends to zero, giving the asserted sign
limit.

For $0<\gamma<1$, couple
\[
 Z_N=N^{-\gamma/2}Q_N^{1/2}\xi,\qquad
 Y_N=N^{-\gamma/2}K^{1/2}\xi,
\]
where $\xi\sim\mathcal N(0,I_M)$ is independent of $Q_N$.
The square-root map is locally Lipschitz near the positive-definite
$K$. For fixed $\delta,\epsilon>0$, this gives
\[
 \Prb\{\|Z_N-Y_N\|>\delta\}
 \le C_\epsilon e^{-c_\epsilon N}
   +\Prb\left\{\|\xi\|>
             \frac{\delta N^{\gamma/2}}{C\epsilon}\right\}.
\]
Divide logarithms by $N^\gamma$, let $N\to\infty$, and then let
$\epsilon\downarrow0$. The upper limit is $-\infty$, since
$\gamma<1$. This is exponential equivalence at speed $N^\gamma$.
The Gaussian family $Y_N$ has the claimed quadratic rate, and
exponential equivalence transfers its large-deviation principle.
\end{proof}

\begin{proposition}[Exact normalized parity rate]
\label{prop:rate}
For $0<\gamma<1$,
\begin{equation}\label{eq:exact-rate}
 \mathcal C_{\gamma,a}(f_d)
 =\frac{a^2}{2}\left(\frac1{\lambda_d}-\frac1{\lambda_0}\right).
\end{equation}
\end{proposition}

\begin{proof}
For a Boolean sign function $g$, define
\[
 E_K(g)=\min\left\{\frac12z^TK^{-1}z:
                           g(x)z(x)\ge a\text{ for all }x\right\}.
\]
The feasible set is closed and the objective strictly convex and
coercive, so it has a unique minimizer. Scaling that minimizer by
$1+\eta$ puts it in the strict margin orthant and approaches its
energy as $\eta\downarrow0$. The open lower and closed upper bounds
in Proposition~\ref{prop:kernel-limit} therefore agree:
\[
 \lim_{N\to\infty}-N^{-\gamma}\log\Prb(A_g)=E_K(g).
\]
The global event is the disjoint union of finitely many such
orthants. Its rate is $\min_g E_K(g)$, so the normalized prior rate
is $E_K(g)-\min_hE_K(h)$.

If a Boolean vector $g$ has normalized eigenvalue $\lambda$, then
$z=ag$ minimizes its margin energy. Its gradient is
$K^{-1}(ag)=ag/(M\lambda)$, whose inner product with every
feasible displacement is nonnegative. Hence
$E_K(g)=a^2/(2\lambda)$.

For any full-margin vector $z$,
\[
 \frac12z^TK^{-1}z
 \ge \frac{\|z\|_2^2}{2M\lambda_0}
 \ge \frac{a^2}{2\lambda_0}.
\]
The constant sign attains equality. Full parity has eigenvalue
$\lambda_d$, proving~\eqref{eq:exact-rate}.
\end{proof}

For $\gamma\le0$, the limiting sign probabilities are positive,
so negative logarithms converge as in~\eqref{eq:complexity}.
Applying Lemma~\ref{lem:gaussian-sign} to full parity and the
appropriate margin proves~\eqref{eq:nonpositive-bound}.
Combining Proposition~\ref{prop:rate} with
Lemma~\ref{lem:spectrum} proves~\eqref{eq:positive-bound}.

\section{An explicit small circuit for parity}

\begin{proposition}[Alternating Hamming-weight thresholds]
\label{prop:circuit}
For every $d\ge1$ and $L\ge2$, there is a depth-$L$ $\tanh$ network
of width $d+1$, using only the appended input bias, whose output
approximates $f_d$ uniformly within $1/8$ and whose weights have
magnitude at most $(d+1)\log(16(d+1))$.
\end{proposition}

\begin{proof}
Set $A=\log(16(d+1))$, $s(x)=\sum_i x_i$, and for $1\le j\le d$
put
\[
 t_j=-d+2j-1,\qquad q_j(x)=\sign(s(x)-t_j),\qquad
 c_j=(-1)^{d-j}.
\]
Let $\alpha_d=1$ for even $d$ and $\alpha_d=0$ for odd $d$.
The exact identity
\begin{equation}\label{eq:parity-thresholds}
 f_d(x)=\alpha_d+\sum_{j=1}^d c_jq_j(x)
\end{equation}
follows by writing $r$ for the number of $+1$ bits. Then $q_j=+1$
exactly when $j\le r$. At $r=0$ both sides equal $(-1)^d$, and
each increment of $r$ changes both sides by $2(-1)^{d-r}$.

The first hidden layer contains
$h_j^1(x)=\tanh(A(s(x)-t_j))$ for $1\le j\le d$, together with
the constant neuron $h_0^1(x)=\tanh A$. Each $s(x)-t_j$ is a nonzero
odd integer. All these activations consequently have their desired
sign and differ from it by at most
$\delta=1/[8(d+1)]$.

For every additional hidden layer use diagonal weights $A$ and
zero off-diagonal weights. If an activation has the correct sign
and magnitude at least $1/2$, its successor differs from that
sign by at most $2e^{-A}=\delta$. Thus the same approximation
holds at every hidden layer, without accumulated error.

Use the linear readout
\[
 z(x)=\alpha_d h_0^{L-1}(x)
             +\sum_{j=1}^d c_jh_j^{L-1}(x).
\]
By~\eqref{eq:parity-thresholds}, its uniform error is at most
$(d+1)\delta=1/8$. In particular its correct margin is at least
$7/8$. The largest weight is at most $A\max(1,d-1)$, including
the first-layer threshold biases. This is bounded by $(d+1)A$.
\end{proof}

This proves~\eqref{eq:circuit-bound} and completes the proof of
Theorem~\ref{thm:main}. If a circuit implementing an arbitrary
fixed margin $a$ is desired, rescale the readout by a constant
depending only on $a$; the polynomial size conclusion is unchanged.

\section{What remains a learning question}\label{sec:learning}

The theorem separates circuit complexity from prior complexity
throughout $\gamma<1$, including $k=d$ for a $k$-bit XOR. It does
not establish that the actual posterior requires exponentially many
examples to achieve small test error or correct-label probability
strictly greater than $2/3$. A low prior probability can be overcome
by a small sample if that sample eliminates the competing functions.

For example, consider the Gaussian process
$G(x)=A+Bf(x)$, where $A$ and $B$ are independent nondegenerate
centered Gaussians and $f$ is a nonconstant parity. Even after any
global margin conditioning, its possible sign functions are only
$1,-1,f,-f$. Once examples with both target labels have been seen,
$f$ is the unique compatible sign function and every prediction
has correct-label probability one. The variance of $B$ can be
arbitrarily small, making the initial probability of $f$ arbitrarily
unfavorable without changing that conclusion. This example does
not contradict Lemma~\ref{lem:gaussian-sign}: its covariance is
singular, and $f$ is not a minimum-eigenvalue eigenvector.

For the networks studied here, the same kernel limit identifies the
posterior completely at $\gamma\le0$ and its leading energy at
$0<\gamma<1$. At $\gamma\le0$ it is a Gaussian sign
classifier with the margin specified in
Proposition~\ref{prop:kernel-limit}. At $0<\gamma<1$, conditioned
on observed training signs, its mass concentrates on compatible
sign functions minimizing $E_K$. The optimization retains
$|z(x)|\ge a$ at every unseen input. Its leading optimization uses
a fixed kernel and a nonconvex union of full margin orthants as its
feasible set. This statement does not determine posterior weights
between tied minimizers. In fact, the companion note on Gaussian
conditioning gives an actual-network example where non-Gaussian
corrections change those weights and a $2/3$-confidence prediction
for $1/2\le\gamma<1$. Ordinary Gaussian-process regression or a hard-margin support
vector machine using only the observed inputs is a different rule.

An exponential sample lower bound for this precise classifier
requires control of compatible competitors and posterior marginal
sign probabilities, beyond the prior separation proved here.
The confidence criterion $>2/3$ makes that distinction especially
important: a tiny posterior mass on the exact target function can
coexist with highly confident correct marginal predictions at
every input.

\end{document}
